\section{Cơ chế phòng thủ \ours}\label{sec:defence}

\ours{} dựa trên ba nguyên tắc. Thứ nhất, mô hình không phải ranh giới an ninh, vì một hành động chỉ có hiệu lực khi
broker thực thi nó. Theo cách gọi kinh điển, broker đóng vai trò một reference monitor tại các sink. Thứ hai, người
dùng là chủ thể đáng tin duy nhất. Quyền được suy ra từ nguyên văn $u$, không từ nội dung công cụ và không từ việc một
mô hình diễn giải ý định. Thứ ba, việc chặn không được làm dừng nhiệm vụ. Hành động bị từ chối được báo lại để tác tử
hoàn thành phần việc hợp lệ. Hình~\ref{fig:pipeline} mô tả pipeline. Bảng~\ref{tab:versions} so sánh \ours{} với cấu
hình chỉ có luật \cRules, dùng để tách đóng góp của từng lớp. Tập phát triển Dev và hai tập kiểm thử niêm phong \tOne,
\tTwo{} được mô tả ở Mục~\ref{sec:protocol}.

% Hình 1, vẽ bằng TikZ (Overleaf tự vẽ, không render trên máy). Thay fig1_pipeline.pdf (bản đó vẽ BG-v2).
% Bố cục chép từ images/src/fig1_pipeline.py, đơn vị cũ nhân 1,2 cm. Tổng rộng 12,0 cm < \textwidth 12,2 cm.
% Nhãn tiếng Anh cho bản nộp (thay từng nhãn tương ứng):
%   User / instruction $u$ / (trusted) | LLM agent / (ReAct) | Tool Broker (host side)
%   Schema and limits · trust and sensitivity labels · confidential values $C$
%   before every sink call | R-Dest: destination must appear in $u$
%   R-Flow: confidential value leaves only if written in $u$
%   R-Signal: after a flagged observation, deny sinks unless $u$ allows the call
%   at the answer · on observations · on denial | R-Flow + R-Redact: redact each confidential value in the answer
%   Sandwich: repeat $u$ after each observation | Denial feedback: tool not executed, agent continues
%   Read tools: untrusted output, may carry injections | Sinks: e-mail, webhook, mocks, no network
%   redacted answer | deterministic check | model-dependent
\begin{figure}[!htbp]
\centering
\begin{tikzpicture}[
  font=\fontsize{5.6}{6.6}\selectfont,
  box/.style={rounded corners=1pt, line width=0.5pt, align=center, inner sep=1pt},
  rule/.style={box, draw=bgblue, fill=bgblue!12},
  soft/.style={box, draw=bgorange, fill=bgorange!10, dash pattern=on 2pt off 1.3pt},
  neutral/.style={box, draw=black!60, fill=black!5},
  arr/.style={-{Stealth[length=3pt]}, line width=0.5pt, draw=black!65},
  lab/.style={font=\fontsize{6.4}{7}\selectfont},
  note/.style={font=\fontsize{5.4}{6.4}\selectfont\itshape, text=black!65, anchor=west},
]
% cột trái: người dùng và tác tử
\node[neutral, minimum width=1.5cm, minimum height=1.14cm, text width=1.4cm] at (0.75,4.23)
  {Người dùng\\ chỉ thị $u$\\ (đáng tin)};
\node[neutral, minimum width=1.5cm, minimum height=1.56cm, text width=1.4cm,
  font=\fontsize{6.6}{7.6}\selectfont\bfseries] at (0.75,1.98) {Tác tử LLM\\ (ReAct)};
\draw[arr] (0.744,3.66) -- (0.744,2.76);

% khung broker
\draw[rounded corners=2pt, line width=0.6pt] (2.22,0.66) rectangle (9.12,5.46);
\node[font=\fontsize{7}{8}\selectfont\bfseries] at (5.67,5.22) {Tool Broker (phía host)};
\node[rule, minimum width=6.612cm, minimum height=0.42cm, text width=6.5cm] at (5.67,4.674)
  {Kiểm schema và giới hạn $\cdot$ gán nhãn tin cậy, nhạy cảm $\cdot$ trích giá trị mật $C$};

\node[note] at (2.33,4.236) {trước mỗi lời gọi sink};
\node[rule, minimum width=2.124cm, minimum height=0.96cm, text width=2.0cm] at (3.426,3.6)
  {\textsf{R-Dest}\\ đích phải có trong $u$};
\node[rule, minimum width=2.124cm, minimum height=0.96cm, text width=2.0cm] at (5.670,3.6)
  {\textsf{R-Flow}\\ giá trị mật chỉ được ra\\ nếu có trong $u$};
\node[rule, minimum width=2.124cm, minimum height=0.96cm, text width=2.0cm] at (7.914,3.6)
  {\textsf{R-Signal}\\ có quan sát bị gắn cờ:\\ chặn sink trừ khi $u$ cho phép};

\node[note] at (2.33,2.88) {ở câu trả lời $\cdot$ trên đường quan sát $\cdot$ khi từ chối};
\node[rule, minimum width=2.124cm, minimum height=0.96cm, text width=2.0cm] at (3.426,2.22)
  {\textsf{R-Flow} + \textsf{R-Redact}\\ che từng giá trị mật\\ trong câu trả lời};
\node[soft, minimum width=2.124cm, minimum height=0.96cm, text width=2.0cm] at (5.670,2.22)
  {Sandwich\\ nhắc lại $u$ sau\\ mỗi quan sát};
\node[rule, minimum width=2.124cm, minimum height=0.96cm, text width=2.0cm] at (7.914,2.22)
  {Phản hồi từ chối\\ không chạy công cụ,\\ tác tử làm tiếp};

% chú giải
\node[rule, minimum width=0.336cm, minimum height=0.24cm] at (2.568,1.056) {};
\node[note, font=\fontsize{5.4}{6.4}\selectfont] at (2.80,1.056) {kiểm tất định};
\node[soft, minimum width=0.336cm, minimum height=0.24cm] at (4.968,1.056) {};
\node[note, font=\fontsize{5.4}{6.4}\selectfont] at (5.20,1.056) {phụ thuộc mô hình};

% tác tử <-> broker
\draw[arr] (1.5,2.46) -- node[lab, above=-1pt] {$a_t$} (2.22,2.46);
\draw[arr] (2.22,1.98) -- node[lab, above=-1pt] {$o_t$} (1.5,1.98);
\draw[arr] (1.5,1.44) -- node[lab, above=-1pt] {$y$} (2.22,1.44);

% cột công cụ
\node[neutral, minimum width=2.16cm, minimum height=1.14cm, text width=2.06cm] at (10.92,3.69)
  {Công cụ đọc\\ đầu ra không đáng tin,\\ có thể chứa lệnh chèn};
\node[neutral, minimum width=2.16cm, minimum height=0.96cm, text width=2.06cm] at (10.92,2.22)
  {Sink: e-mail, webhook\\ bản giả, không ra mạng};
\draw[arr] (9.12,3.90) -- node[lab, above=-1pt] {$a_t$} (9.84,3.90);
\draw[arr] (9.84,3.42) -- node[lab, above=-1pt] {$o_t$} (9.12,3.42);
\draw[arr] (9.12,2.22) -- node[lab, above=-1pt] {$a_t$} (9.84,2.22);

% câu trả lời đã che rời broker
\draw[arr] (9.12,1.02) -- (10.26,1.02);
\node[anchor=west, text=black!65] at (10.30,1.02) {câu trả lời đã che};
\end{tikzpicture}
\caption{Pipeline của \ours. Mọi hành động $a_t$ và câu trả lời $y$ đều đi qua Tool Broker. Ô liền nét là kiểm tra
tất định phía host, ô nét đứt là lớp phụ thuộc mô hình.}\label{fig:pipeline}
\end{figure}

% Thành phần từng cấu hình (đối chiếu configs/final/final_config_v3.json: D2+P2 = D2 runtime --levels A1 + P2).
% 13/11: bỏ BG-v2 (D3DEST) khỏi bài; BrokerGuard = D2+P2.
\begin{table}[!htbp]
\centering
\caption{Thành phần của \ours{} và cấu hình chỉ có luật. \checkmark: có, -: không.}\label{tab:versions}
\footnotesize
\begin{adjustbox}{max width=\textwidth}
\begin{tabular}{@{}lccccccl@{}}
\toprule
Cấu hình & \rDest & \rFlow & \rRedact & \rSig & Sandwich & Phản hồi & Đánh giá trên \\
\midrule
\cRules{} (chỉ luật) & \checkmark & \checkmark & \checkmark & \checkmark & - & \checkmark & Dev, \tOne{} (tách lớp) \\
\textbf{\ours} & \checkmark & \checkmark & \checkmark & \checkmark & \checkmark & \checkmark & Dev, \tOne, \tTwo \\
\bottomrule
\end{tabular}
\end{adjustbox}
\end{table}


\paragraph{Broker và nhãn nguồn gốc.} Khi không có phòng thủ (\cNone), broker chỉ kiểm tra schema và giới hạn của công
cụ. Với \ours, broker lưu thêm nhãn của từng quan sát và tập giá trị $C$ trích từ các quan sát mật. Việc trích giá trị
dựa trên quy tắc từ vựng, không phụ thuộc ngôn ngữ, gồm e-mail, URL, định danh từ sáu ký tự trở lên và số dài. Thao
tác đọc không bao giờ bị chặn. Mọi phép so sánh dùng hàm chuẩn hóa $\norm(\cdot)$, gồm chuẩn hóa Unicode NFKC, bỏ ký tự
định dạng như ký tự độ rộng bằng không, và casefold. Ta viết $x \sub z$ khi $\norm(x)$ là một chuỗi con của
$\norm(z)$. Đây là một quan hệ được kiểm tra trực tiếp, không dựng tập mọi chuỗi con.

\paragraph{Luật \rDest{} (cấp quyền đích).} Một lời gọi sink $a$ chỉ được thực thi nếu
\begin{equation}\label{eq:dest}
\dest(a) \sub u .
\end{equation}
Luật này không phân tích ý định và không dùng danh sách cụm từ. Vì vậy, một đoạn chèn \mCR{} tránh động từ ``gửi'' hay
tách nhỏ địa chỉ vẫn thất bại vì cùng lý do với \mER{}: đích mà nó cần không có trong $u$. Luật dựa trên một giả định.
Mọi đích xuất hiện trong $u$ được coi là đã được cấp quyền. Giả định này sai theo hai hướng. Một đích chỉ xuất hiện
trong câu cấm vẫn được tính là đã cấp quyền, và cơ chế \mUS{} được thiết kế để đo đúng trường hợp này. Ngược lại, yêu
cầu nhắc tới đích một cách gián tiếp (``địa chỉ trong hợp đồng'') không được đáp ứng. Ngoài ra, so khớp chuỗi con chấp
nhận cả một phần của đích người dùng viết, chẳng hạn phần trước tên miền của một địa chỉ e-mail. Broker kiểm tra dạng
chuẩn hóa nhưng công cụ nhận đích nguyên dạng. Trong benchmark các sink là bản giả nên hai điểm này không gây hại thật.
Với công cụ thật, nên so khớp trên đích đã phân tích cấu trúc và thực thi đúng dạng đã kiểm.

\paragraph{Luật \rFlow{} và \rRedact{} (luồng giá trị mật).} Với một lời gọi sink $a$, broker từ chối nếu
\begin{equation}\label{eq:flow}
\exists\, v \in C:\; v \sub \pay(a) \;\land\; \lnot\,(v \sub u) .
\end{equation}
Với câu trả lời cuối, điều kiện tương tự được áp dụng cho $y$ thay cho $\pay(a)$. Khi đó \rRedact{} không loại bỏ cả
câu trả lời mà thay từng giá trị vi phạm bằng \code{[REDACTED]}. Việc coi $y$ là một sink đóng kênh \mAL{}, vốn bị các
bộ lọc lời gọi công cụ bỏ sót. Phạm vi bảo đảm của luật này hẹp hơn một chính sách kiểm soát luồng thông tin. Luật chỉ
cưỡng chế tất định các luồng \emph{nguyên văn}. Một giá trị bị mô hình biến đổi, chẳng hạn được định dạng lại, tách ra
hay mã hóa, không còn thỏa $v \sub \pay(a)$ và sẽ đi qua. Luật cũng không đóng các kênh ngầm, ví dụ việc chọn gửi hay
không gửi. Giá trị mật không có dạng mà bộ trích nhận ra cũng không được bảo vệ. Chúng tôi chưa đo độ phủ của bộ trích
so với tập giá trị mật mà bộ chấm dùng.

\paragraph{Luật tín hiệu chèn lệnh \rSig.} Hai luật trên không bao quát việc thao túng một sink mà người dùng
\emph{đã} nêu bằng giá trị không mật, tức cơ chế \mPT. Một bộ phát hiện cụm từ gắn cờ các quan sát chứa văn bản giống
chỉ thị. Gọi $F$ là tập quan sát bị gắn cờ và $\mathrm{Gram}_u(a)$ là điều kiện một văn phạm chặt trên $u$ cho phép
lời gọi $a$. Broker từ chối $a$ nếu
\begin{equation}\label{eq:signal}
F \neq \emptyset \;\land\; \lnot\,\mathrm{Gram}_u(a) .
\end{equation}
Trên Dev, bỏ luật này làm số họ gây hại tăng thêm 2 đến 4 ở mỗi mô hình nền. Đổi lại, đây là nguồn chặn nhầm lớn nhất.
Trên \tOne, dưới \cRules, luật này chặn nhầm 22 ca benign\_plain, trong khi \rDest{} chỉ chặn nhầm 3 ca. Bộ phát hiện
dùng danh sách cụm từ nên phụ thuộc ngôn ngữ (Mục~\ref{sec:discussion}).

\paragraph{Lời nhắc và phản hồi từ chối.} Sau mỗi quan sát, $u$ được nhắc lại cùng một câu yêu cầu bỏ qua mọi chỉ thị
nằm trong dữ liệu công cụ. Cách làm này giống phòng thủ sandwich và \code{repeat\_user\_prompt} của AgentDojo. Đây là
lớp duy nhất phụ thuộc vào mô hình, và nó bù cho các luật ở những chỗ luật không thấy được (Mục~\ref{sec:results}).
Khi một luật từ chối, công cụ không được chạy và tác tử nhận một quan sát có cấu trúc báo rằng hành động đã bị từ
chối. Nhờ đó tác tử vẫn hoàn thành được các bước còn lại. Phản hồi từ chối có mặt ở mọi cấu hình có luật nên chưa được
tách riêng.

\paragraph{Chọn cấu hình và chi phí.} Mọi lựa chọn được thực hiện trên Dev trong các vòng chỉnh khai báo trước. \ours{}
được tạo ở vòng chỉnh thứ hai bằng cách ghép lớp luật \cRules{} với \cSand, hai thành phần đã có, không đổi tham số
nào. Cấu hình này được đóng băng
trước khi chạy \tOne{} và không thay đổi về sau. \ours{} không gọi thêm mô hình nào. Mỗi kiểm tra là một phép tìm chuỗi
con trên $u$, $\pay(a)$ hoặc $y$, rẻ hơn nhiều so với một bước sinh của LLM. Chúng tôi chưa đo độ trễ. Lời nhắc làm
ngữ cảnh dài thêm một bản $u$ sau mỗi quan sát.
